\section{Conclusion and Next Steps}

The study leaves us with three observations that together frame the
design space for a multi-tenant vector store.

\mypar{Search is per segment in parallel.} Vectors are arranged
into sealed segments, and every query fans out to all of them,
scans each local index, and merges the partial results into a
global top-$k$. Per-query cost scales with both the number of
segments and the size of each segment's index, so any sharing
strategy must either reduce the number of segments a query
consults or amortize per-segment cost across tenants.

\mypar{Same data does not imply same segments.} Identical
logical data loaded into two collections produces different
segment cuts in stock Milvus, because flush jitter and background
compaction shuffle vectors across segment boundaries during
ingestion. Disabling both sources of non-determinism collapses
the per-segment overlap matrix to a diagonal, but the index and
data files within each segment still diverge across loads.
Vector-level overlap is therefore necessary but not sufficient
for sharing storage; a layout that makes segments byte-comparable
across collections by construction is required before sharing
can be exploited end to end.

\mypar{Sharing a collection has a real cost.} Co-locating
tenants in a single collection through a shared partition
removes the storage duplication of the fully replicated layout,
but query throughput under concurrent traffic collapses by an
order of magnitude relative to the solo case. The bottleneck is
not CPU but the page cache: the combined working set overflows
RAM and queries refault pages that concurrent queries have just
evicted, so the system enters a thrashing regime in which adding
threads reduces throughput. A workable shared layout must
therefore be paired with a memory allocation strategy that
bounds cross-tenant interference rather than relying on the
operating system page cache.

\mypar{Next steps.} These observations point to three follow-up
directions:
\begin{enumerate*}[label=(\arabic*)]
    \item \textbf{Make the search path content aware,} so that a
    query is routed only to the segments that hold its candidate
    top-$k$ neighbors rather than fanned out to every segment,
    and quantify how much per-segment HNSW~\cite{malkov2018hnsw}
    work could additionally be reused across tenants whose
    queries hit the same segments;
    \item \textbf{align segments by construction,} replacing
    implicit ingestion-time alignment with an explicit
    content-addressed segment layout, so that two collections
    with overlapping data share the same segment files rather
    than two independently flushed copies;
    \item \textbf{bound cross-tenant memory contention,}
    developing a memory allocation policy for shared-partition
    collections that reserves working-set budget per tenant or
    per partition, and extending the evaluation beyond two
    tenants.
\end{enumerate*}
